\section{Charm-quark self-scale \texorpdfstring{$\overline{\rm MS}$}{MSbar} transport closeout}

The APF up-family trace theorem gives
\begin{equation}
  m_c^{\rm APF\text{-}TRACE}=0.680091224926\ {\rm GeV}.
\end{equation}
This value is an internal APF trace anchor. It is not, by definition, a pole,
Monte-Carlo, threshold, lattice, or \(\overline{\rm MS}\) mass.

The natural short-distance comparison target for charm is the self-scale
running mass
\begin{equation}
  m_c(m_c)_{\overline{\rm MS}}.
\end{equation}
The PDG 2025 listing gives
\begin{equation}
  m_c(m_c)_{\overline{\rm MS}}=1.2730\pm0.0046\ {\rm GeV}\quad({\rm CL}=90\%).
\end{equation}
The direct residual is therefore
\begin{equation}
  m_c^{\rm APF\text{-}TRACE}-m_c(m_c)_{\overline{\rm MS}}
  =-0.592908775074\ {\rm GeV},
\end{equation}
namely about \(-46.6\%\). This rejects a direct self-scale
\(\overline{\rm MS}\) export.

The correct status is consequently
\begin{equation}
\boxed{
  m_c^{\rm APF\text{-}TRACE}:
  [P_{\rm trace}+P_{\overline{\rm MS}\ route\ knockout}]
}
\end{equation}
not
\begin{equation}
\boxed{
  m_c^{\rm APF\to\overline{MS}}(m_c):[P_{\rm export\ candidate}].
}
\end{equation}
The first failed gate is
\begin{equation}
\boxed{\texttt{APF\_CHARM\_TRACE\_TO\_MSBAR\_NORMALIZATION\_OR\_TRANSPORT\_MAP}.}
\end{equation}
No physical charm mass, pole mass, lattice mass, or PDG value is used as an input
in the APF trace computation; the comparison is external validation only.
